% --- Archon physics formalization source begin ---
% archon:physics
% archon:covers PhyXMiniProblems/problem_phyx_mini_0136.lean
% archon:source-report reports/phyx_mini/problem_phyx_mini_0136.source.json

\chapter{Physics problem phyx\_mini\_0136}
\label{ch:PhyXMiniProblems_problem_phyx_mini_0136}

\paragraph{Problem source.}
\#\# Physical scenario

In a film projector, the film acts as the object whose image is projected on a screen (figure).A \textbackslash{}(105\textbackslash{},\textbackslash{}mathrm\{mm\}\textbackslash{})-focal-length lens is to project an image on a screen \textbackslash{}(25.5\textbackslash{},\textbackslash{}mathrm\{m\}\textbackslash{}) away.The film is \textbackslash{}(24\textbackslash{},\textbackslash{}mathrm\{mm\}\textbackslash{}) wide.

\#\# Question

How wide will the picture be on the screen?

\#\# Answer choices

- A: A: 3.8m

- B: B: 4.8m

- C: C: 5.8m

- D: D: 6.8m

\#\# Figure caption (auxiliary; use the image as primary evidence)

The image depicts a diagram illustrating the basic concept of projection using a lens system. Key elements include:

- **Film**: Positioned on the left, represented by a vertical gray line. It is the starting point for light rays.
- **Lens**: Centrally located, depicted as a vertical oval shape with a light green fill, symbolizing its curvature.
- **Light Rays**: Shown as magenta lines with arrows indicating the direction of light traveling from the film through the lens and converging onto a screen.
- **Screen**: Positioned on the right, represented by a vertical gray line, where the light rays converge.

The light rays start from the film, pass through the lens, and focus on the screen, demonstrating the process of image projection. The text "Film," "Lens," and "Screen" labels each corresponding part of the diagram.

\#\# Dataset metadata

Geometrical Optics; Physical Model Grounding Reasoning, Spatial Relation Reasoning

\paragraph{Recorded answer/context.}
C: C: 5.8m

\paragraph{Figure/image path.}
/root/proposal\_for\_physic/science-mango/phyx\_mini\_run/phyx\_data/test\_image/136.png

\paragraph{Formalization target.}
create a compiling Lean file with sorry bodies at `PhyXMiniProblems/problem\_phyx\_mini\_0136.lean`. The Lean declarations must preserve the physical quantities, dimensions or dimensional roles, figure labels, governing-law hypotheses, and final relation expressed by this problem.
Use Mathlib/Physlib names found through LeanExplore where available. If a physics API is missing, introduce faithful local abstractions rather than scalar placeholder aliases.

\begin{theorem}[Physics formalization target]
\label{thm:physics:phyx_mini_0136:target}
\lean{PhyXMiniProblems.ProblemPhyXMini0136.problem_phyx_mini_0136}
\uses{def:physics:phyx-mini-0136:phyxminiproblems-problemphyxmini0136-metersvalue, def:physics:phyx-mini-0136:phyxminiproblems-problemphyxmini0136-filmprojectorsetup, def:physics:phyx-mini-0136:phyxminiproblems-problemphyxmini0136-hasstatedprojectorreadouts, def:physics:phyx-mini-0136:phyxminiproblems-problemphyxmini0136-matchesprojectorfigure, def:physics:phyx-mini-0136:phyxminiproblems-problemphyxmini0136-hasphysicalprojectorlengths, def:physics:phyx-mini-0136:phyxminiproblems-problemphyxmini0136-obeysthinlensequation, def:physics:phyx-mini-0136:phyxminiproblems-problemphyxmini0136-obeysparaxialwidthmagnificationlaw, def:physics:phyx-mini-0136:phyxminiproblems-problemphyxmini0136-answerchoicematchespicture}
The assigned autoformalize agent should translate this physics problem into Lean declarations in the covered file, with theorem and lemma proof bodies written as `by sorry`.
\end{theorem}
\begin{proof}
This is an autoformalization task, not a proof task. Produce statements that can later be proved by the physics prover without weakening the physical model.
\end{proof}

% --- Archon declaration topology begin ---
\section{Declaration topology}
\begin{definition}[Optical Length]
  \label{def:physics:phyx-mini-0136:phyxminiproblems-problemphyxmini0136-opticallength}
  \lean{PhyXMiniProblems.ProblemPhyXMini0136.OpticalLength}
  A signed physical length whose numerical value changes coherently with unit choice
\end{definition}
\begin{proof}
  This is the declared model definition of Optical Length.
\end{proof}
\begin{definition}[choices With Length Unit]
  \label{def:physics:phyx-mini-0136:phyxminiproblems-problemphyxmini0136-choiceswithlengthunit}
  \lean{PhyXMiniProblems.ProblemPhyXMini0136.choicesWithLengthUnit}
  Unit choices obtained from SI by changing only the length unit
\end{definition}
\begin{proof}
  This is the declared model definition of choices With Length Unit.
\end{proof}
\begin{definition}[length Readout]
  \label{def:physics:phyx-mini-0136:phyxminiproblems-problemphyxmini0136-lengthreadout}
  \lean{PhyXMiniProblems.ProblemPhyXMini0136.lengthReadout}
  \uses{def:physics:phyx-mini-0136:phyxminiproblems-problemphyxmini0136-opticallength, def:physics:phyx-mini-0136:phyxminiproblems-problemphyxmini0136-choiceswithlengthunit}
  Scalar readout of a physical length in the selected length unit
\end{definition}
\begin{proof}
  This is the declared model definition of length Readout.
\end{proof}
\begin{definition}[meters Value]
  \label{def:physics:phyx-mini-0136:phyxminiproblems-problemphyxmini0136-metersvalue}
  \lean{PhyXMiniProblems.ProblemPhyXMini0136.metersValue}
  \uses{def:physics:phyx-mini-0136:phyxminiproblems-problemphyxmini0136-opticallength, def:physics:phyx-mini-0136:phyxminiproblems-problemphyxmini0136-lengthreadout}
  Scalar readout of a physical length in metres
\end{definition}
\begin{proof}
  This is the declared model definition of meters Value.
\end{proof}
\begin{definition}[millimeters Value]
  \label{def:physics:phyx-mini-0136:phyxminiproblems-problemphyxmini0136-millimetersvalue}
  \lean{PhyXMiniProblems.ProblemPhyXMini0136.millimetersValue}
  \uses{def:physics:phyx-mini-0136:phyxminiproblems-problemphyxmini0136-opticallength, def:physics:phyx-mini-0136:phyxminiproblems-problemphyxmini0136-lengthreadout}
  Scalar readout of a physical length in millimetres
\end{definition}
\begin{proof}
  This is the declared model definition of millimeters Value.
\end{proof}
\begin{definition}[Thin Lens Kind]
  \label{def:physics:phyx-mini-0136:phyxminiproblems-problemphyxmini0136-thinlenskind}
  \lean{PhyXMiniProblems.ProblemPhyXMini0136.ThinLensKind}
  The optical type of the lens drawn in the projector figure
\end{definition}
\begin{proof}
  This is the declared model definition of Thin Lens Kind.
\end{proof}
\begin{definition}[Image Nature]
  \label{def:physics:phyx-mini-0136:phyxminiproblems-problemphyxmini0136-imagenature}
  \lean{PhyXMiniProblems.ProblemPhyXMini0136.ImageNature}
  Whether the transmitted rays physically meet at the image
\end{definition}
\begin{proof}
  This is the declared model definition of Image Nature.
\end{proof}
\begin{definition}[Projector Element]
  \label{def:physics:phyx-mini-0136:phyxminiproblems-problemphyxmini0136-projectorelement}
  \lean{PhyXMiniProblems.ProblemPhyXMini0136.ProjectorElement}
  The three labeled components ordered along the principal optical axis
\end{definition}
\begin{proof}
  This is the declared model definition of Projector Element.
\end{proof}
\begin{definition}[Figure Ray]
  \label{def:physics:phyx-mini-0136:phyxminiproblems-problemphyxmini0136-figureray}
  \lean{PhyXMiniProblems.ProblemPhyXMini0136.FigureRay}
  The three representative magenta ray paths visible in the source figure
\end{definition}
\begin{proof}
  This is the declared model definition of Figure Ray.
\end{proof}
\begin{definition}[Principal Axis Orientation]
  \label{def:physics:phyx-mini-0136:phyxminiproblems-problemphyxmini0136-principalaxisorientation}
  \lean{PhyXMiniProblems.ProblemPhyXMini0136.PrincipalAxisOrientation}
  Orientation of the principal optical axis shown in the figure
\end{definition}
\begin{proof}
  This is the declared model definition of Principal Axis Orientation.
\end{proof}
\begin{definition}[Propagation Direction]
  \label{def:physics:phyx-mini-0136:phyxminiproblems-problemphyxmini0136-propagationdirection}
  \lean{PhyXMiniProblems.ProblemPhyXMini0136.PropagationDirection}
  Direction in which light propagates through the projector
\end{definition}
\begin{proof}
  This is the declared model definition of Propagation Direction.
\end{proof}
\begin{definition}[Thin Lens]
  \label{def:physics:phyx-mini-0136:phyxminiproblems-problemphyxmini0136-thinlens}
  \lean{PhyXMiniProblems.ProblemPhyXMini0136.ThinLens}
  \uses{def:physics:phyx-mini-0136:phyxminiproblems-problemphyxmini0136-opticallength, def:physics:phyx-mini-0136:phyxminiproblems-problemphyxmini0136-thinlenskind}
  A thin lens together with its physical focal length and qualitative kind
\end{definition}
\begin{proof}
  This is the declared model definition of Thin Lens.
\end{proof}
\begin{definition}[Film Projector Setup]
  \label{def:physics:phyx-mini-0136:phyxminiproblems-problemphyxmini0136-filmprojectorsetup}
  \lean{PhyXMiniProblems.ProblemPhyXMini0136.FilmProjectorSetup}
  \uses{def:physics:phyx-mini-0136:phyxminiproblems-problemphyxmini0136-opticallength, def:physics:phyx-mini-0136:phyxminiproblems-problemphyxmini0136-imagenature, def:physics:phyx-mini-0136:phyxminiproblems-problemphyxmini0136-projectorelement, def:physics:phyx-mini-0136:phyxminiproblems-problemphyxmini0136-figureray, def:physics:phyx-mini-0136:phyxminiproblems-problemphyxmini0136-principalaxisorientation, def:physics:phyx-mini-0136:phyxminiproblems-problemphyxmini0136-propagationdirection, def:physics:phyx-mini-0136:phyxminiproblems-problemphyxmini0136-thinlens}
  This structure records the Film Projector Setup component of the problem-specific physical model
\end{definition}
\begin{proof}
  This is the declared model definition of Film Projector Setup.
\end{proof}
\begin{definition}[Has Stated Projector Readouts]
  \label{def:physics:phyx-mini-0136:phyxminiproblems-problemphyxmini0136-hasstatedprojectorreadouts}
  \lean{PhyXMiniProblems.ProblemPhyXMini0136.HasStatedProjectorReadouts}
  \uses{def:physics:phyx-mini-0136:phyxminiproblems-problemphyxmini0136-metersvalue, def:physics:phyx-mini-0136:phyxminiproblems-problemphyxmini0136-millimetersvalue, def:physics:phyx-mini-0136:phyxminiproblems-problemphyxmini0136-filmprojectorsetup}
  This structure records the Has Stated Projector Readouts component of the problem-specific physical model
\end{definition}
\begin{proof}
  This is the declared model definition of Has Stated Projector Readouts.
\end{proof}
\begin{definition}[Matches Projector Figure]
  \label{def:physics:phyx-mini-0136:phyxminiproblems-problemphyxmini0136-matchesprojectorfigure}
  \lean{PhyXMiniProblems.ProblemPhyXMini0136.MatchesProjectorFigure}
  \uses{def:physics:phyx-mini-0136:phyxminiproblems-problemphyxmini0136-metersvalue, def:physics:phyx-mini-0136:phyxminiproblems-problemphyxmini0136-figureray, def:physics:phyx-mini-0136:phyxminiproblems-problemphyxmini0136-filmprojectorsetup}
  This structure records the Matches Projector Figure component of the problem-specific physical model
\end{definition}
\begin{proof}
  This is the declared model definition of Matches Projector Figure.
\end{proof}
\begin{definition}[Has Physical Projector Lengths]
  \label{def:physics:phyx-mini-0136:phyxminiproblems-problemphyxmini0136-hasphysicalprojectorlengths}
  \lean{PhyXMiniProblems.ProblemPhyXMini0136.HasPhysicalProjectorLengths}
  \uses{def:physics:phyx-mini-0136:phyxminiproblems-problemphyxmini0136-metersvalue, def:physics:phyx-mini-0136:phyxminiproblems-problemphyxmini0136-filmprojectorsetup}
  Positivity conditions appropriate to physical distance and width magnitudes
\end{definition}
\begin{proof}
  This is the declared model definition of Has Physical Projector Lengths.
\end{proof}
\begin{definition}[Obeys Thin Lens Equation]
  \label{def:physics:phyx-mini-0136:phyxminiproblems-problemphyxmini0136-obeysthinlensequation}
  \lean{PhyXMiniProblems.ProblemPhyXMini0136.ObeysThinLensEquation}
  \uses{def:physics:phyx-mini-0136:phyxminiproblems-problemphyxmini0136-lengthreadout, def:physics:phyx-mini-0136:phyxminiproblems-problemphyxmini0136-filmprojectorsetup}
  This def records the Obeys Thin Lens Equation component of the problem-specific physical model
\end{definition}
\begin{proof}
  This is the declared model definition of Obeys Thin Lens Equation.
\end{proof}
\begin{definition}[Obeys Paraxial Width Magnification Law]
  \label{def:physics:phyx-mini-0136:phyxminiproblems-problemphyxmini0136-obeysparaxialwidthmagnificationlaw}
  \lean{PhyXMiniProblems.ProblemPhyXMini0136.ObeysParaxialWidthMagnificationLaw}
  \uses{def:physics:phyx-mini-0136:phyxminiproblems-problemphyxmini0136-lengthreadout, def:physics:phyx-mini-0136:phyxminiproblems-problemphyxmini0136-filmprojectorsetup}
  This def records the Obeys Paraxial Width Magnification Law component of the problem-specific physical model
\end{definition}
\begin{proof}
  This is the declared model definition of Obeys Paraxial Width Magnification Law.
\end{proof}
\begin{definition}[Answer Choice]
  \label{def:physics:phyx-mini-0136:phyxminiproblems-problemphyxmini0136-answerchoice}
  \lean{PhyXMiniProblems.ProblemPhyXMini0136.AnswerChoice}
  Labels of the four answer choices printed with the problem
\end{definition}
\begin{proof}
  This is the declared model definition of Answer Choice.
\end{proof}
\begin{definition}[picture Width Meters]
  \label{def:physics:phyx-mini-0136:phyxminiproblems-problemphyxmini0136-answerchoice-picturewidthmeters}
  \lean{PhyXMiniProblems.ProblemPhyXMini0136.AnswerChoice.pictureWidthMeters}
  \uses{def:physics:phyx-mini-0136:phyxminiproblems-problemphyxmini0136-answerchoice}
  Picture-width readout in metres printed beside an answer choice
\end{definition}
\begin{proof}
  This is the declared model definition of picture Width Meters.
\end{proof}
\begin{definition}[Is Nearest Tenth Readout]
  \label{def:physics:phyx-mini-0136:phyxminiproblems-problemphyxmini0136-isnearesttenthreadout}
  \lean{PhyXMiniProblems.ProblemPhyXMini0136.IsNearestTenthReadout}
  This def records the Is Nearest Tenth Readout component of the problem-specific physical model
\end{definition}
\begin{proof}
  This is the declared model definition of Is Nearest Tenth Readout.
\end{proof}
\begin{definition}[Answer Choice Matches Picture]
  \label{def:physics:phyx-mini-0136:phyxminiproblems-problemphyxmini0136-answerchoicematchespicture}
  \lean{PhyXMiniProblems.ProblemPhyXMini0136.AnswerChoiceMatchesPicture}
  \uses{def:physics:phyx-mini-0136:phyxminiproblems-problemphyxmini0136-metersvalue, def:physics:phyx-mini-0136:phyxminiproblems-problemphyxmini0136-filmprojectorsetup, def:physics:phyx-mini-0136:phyxminiproblems-problemphyxmini0136-answerchoice, def:physics:phyx-mini-0136:phyxminiproblems-problemphyxmini0136-isnearesttenthreadout}
  A choice agrees to the nearest tenth with the physical picture width
\end{definition}
\begin{proof}
  This is the declared model definition of Answer Choice Matches Picture.
\end{proof}
\begin{lemma}[object Distance In Meters eq three fifty seven over three three eight six]
  \label{lem:physics:phyx-mini-0136:phyxminiproblems-problemphyxmini0136-objectdistanceinmeters-eq-three-fifty-seven-over-three-three-eight-six}
  \lean{PhyXMiniProblems.ProblemPhyXMini0136.objectDistanceInMeters_eq_three_fifty_seven_over_three_three_eight_six}
  \uses{def:physics:phyx-mini-0136:phyxminiproblems-problemphyxmini0136-metersvalue, def:physics:phyx-mini-0136:phyxminiproblems-problemphyxmini0136-filmprojectorsetup, def:physics:phyx-mini-0136:phyxminiproblems-problemphyxmini0136-hasstatedprojectorreadouts, def:physics:phyx-mini-0136:phyxminiproblems-problemphyxmini0136-obeysthinlensequation}
  This lemma records the object Distance In Meters eq three fifty seven over three three eight six component of the problem-specific physical model
\end{lemma}
\begin{proof}
  The conclusion follows from the governing relations and figure readouts named in the statement of object Distance In Meters eq three fifty seven over three three eight six.
\end{proof}
\begin{lemma}[picture Width In Meters eq five zero seven nine over eight seven five]
  \label{lem:physics:phyx-mini-0136:phyxminiproblems-problemphyxmini0136-picturewidthinmeters-eq-five-zero-seven-nine-over-eight-seven-five}
  \lean{PhyXMiniProblems.ProblemPhyXMini0136.pictureWidthInMeters_eq_five_zero_seven_nine_over_eight_seven_five}
  \uses{def:physics:phyx-mini-0136:phyxminiproblems-problemphyxmini0136-metersvalue, def:physics:phyx-mini-0136:phyxminiproblems-problemphyxmini0136-filmprojectorsetup, def:physics:phyx-mini-0136:phyxminiproblems-problemphyxmini0136-hasstatedprojectorreadouts, def:physics:phyx-mini-0136:phyxminiproblems-problemphyxmini0136-hasphysicalprojectorlengths, def:physics:phyx-mini-0136:phyxminiproblems-problemphyxmini0136-obeysthinlensequation, def:physics:phyx-mini-0136:phyxminiproblems-problemphyxmini0136-obeysparaxialwidthmagnificationlaw}
  This lemma records the picture Width In Meters eq five zero seven nine over eight seven five component of the problem-specific physical model
\end{lemma}
\begin{proof}
  The conclusion follows from the governing relations and figure readouts named in the statement of picture Width In Meters eq five zero seven nine over eight seven five.
\end{proof}
% --- Archon declaration topology end ---
% --- Archon physics formalization source end ---
